\section{Finite DAG bounded-domain certificate}
\label{sec:finite-dag-domain-certificate}

Consider a finite acyclic multilinear graph.  For each internal vertex $v$,
let $M_v$ be a declared operator-norm upper bound for its multilinear law and
let $\rho_v(r)$ be a declared upper bound for the normal residual after the
rank-$r$ intermediate projection.  Leaves are restricted to a declared domain
$\|x_\ell\|\le L_\ell$.  The bounds need not be sharp; they are inputs to the
certificate and may be empirical enclosures only if their validation is
recorded separately.

Write $U_v$ for the norm of the projected value enclosure and $D_v$ for the
absolute error enclosure relative to the corresponding unprojected
composition.  For a leaf, $U_\ell=L_\ell$ and $D_\ell=0$.  If the input slots of
$v$ are $v_1,\ldots,v_k$, multilinearity and a telescoping expansion give
\begin{align*}
 U_v &\le M_v\prod_{i=1}^k U_{v_i},\\
 D_v &\le \rho_v(r_v)\prod_{i=1}^k U_{v_i}
       +M_v\sum_{i=1}^k
        \left(\prod_{j\ne i}U_{v_j}\right)D_{v_i}.
\end{align*}
The local term is omitted at an unprojected root.  Crucially, the formula is
indexed by input slots, not by distinct node identifiers: a repeated shared
subexpression is therefore charged once per occurrence, while its stored
value and error bounds are computed only once in topological order.

The recurrence also yields an exact allocation objective for this enclosure
model.  Let $G_{\mathrm{root}}=1$.  For every parent-to-child slot, propagate
\begin{equation*}
 G_{v_i}\mathrel{+}=G_v M_v\prod_{j\ne i}U_{v_j}.
\end{equation*}
Then the root-bound contribution of assigning rank $r$ to a projectable node
$v$ is
\begin{equation*}
 c_v(r)=G_v\rho_v(r)\prod_i U_{v_i}.
\end{equation*}
Since the declared $M_v$ and $U_v$ are rank-independent, the certified root
bound is $\sum_v c_v(r_v)$ (with the root term omitted when the root is not
projected).  Thus an integer budget $\sum_v r_v\le B$ is solved exactly by a
knapsack dynamic program over the finite rank tables.  This is a certificate
and allocator for the declared finite DAG and bounded leaf domain, not an
extremal theorem: sharpness of $M_v$ or $\rho_v(r)$, continuum validity of
empirical estimates, and global optimality outside the supplied enclosure
model remain separate questions.

\subsection{Rank-aware refinement}

The rank-independent value enclosure above is intentionally scalable but can
discard useful information. A fixed rank assignment admits the following
refinement. Let $A_v$ enclose the recursively projected value, let $U_v$
enclose the exact value, and define $B_v=\max\{A_v,U_v\}$. If $L_v(r)$ is a
declared bound for the projected operator $P_v\mu_v$ at rank $r$, then
\begin{align*}
 U_v &\le M_v\prod_i U_{v_i},\\
 A_v &\le L_v(r_v)\prod_i A_{v_i},\\
 D_v &\le \rho_v(r_v)\prod_i A_{v_i}
       +L_v(r_v)\sum_i\left(\prod_{j\ne i}B_{v_j}\right)D_{v_i}.
\end{align*}
The root uses $L_{\mathrm{root}}=M_{\mathrm{root}}$ and has no local normal
term when it is unprojected. The mixed factor $B$ follows by telescoping the
multilinear law one slot at a time, using exact inputs on one side of each
difference and projected inputs on the other. This proves a sound fixed-
assignment certificate for shared DAGs, including repeated slots.

Unlike the rank-independent recurrence, $A_v$ and $B_v$ depend on descendant
ranks. The resulting allocation objective is generally coupled; therefore the
repository's rank-aware optimizer is explicitly an exhaustive small-DAG oracle,
while the earlier separable dynamic program remains the scalable
rank-independent method.
